Fix $k\in\mathbb{N}$. For an integer $n\ge 1$ say that $n$ is \emph{$k$-representable} if there exist
\begin{itemize}
\item a prime $p$, and
\item an integer $t$ with $0\le t\le k$, and
\item exponents $e_1,\dots,e_t\in\mathbb{Z}_{\ge 0}$
\end{itemize}
with
\[
 n = p + \sum_{i=1}^t 2^{e_i}.
\]
(Repeated exponents are allowed a priori.)

Question: Does there exist an absolute $k$ such that every integer $n\ge 2$ is $k$-representable?
(If one literally includes $n=1$, then $1$ is not $k$-representable for any $k$ because $p\ge 2$ and the sum of powers of $2$ is nonnegative.)

Stress points/edge cases:
\begin{itemize}
\item parity: even $n$ force the power-of-two part to be odd when $p$ is odd;
\item the meaning of ``powers of $2$'': allowing repetition vs requiring distinct powers;
\item whether ``every integer'' means all $n\ge 1$ or all sufficiently large $n$.
\end{itemize}

2)
